% GENERATED FILE. Edit canonical lessons, then run npm run generate:books.
% canonical-source-hash: fnv1a64:448dd63fd47b2e15
% canonical-lessons: GU-C488-navdavvu, GU-C488-khavdavvu, GU-C488-paheravvu, GU-C488-ramadvu, GU-C488-pivdavvu

\chapter{Bath Time for the Baby}
\label{ch:gu-bath-time-for-the-baby}
\hlchaptermodality{488}

\begin{chapteropening}
\textbf{By the end of this chapter:} \emph{I can say how I look after a baby.}

Complete the last lesson of chapter 488: I can say how I look after a baby.
\end{chapteropening}

\section[\texorpdfstring{\gu{નવડાવવું} (navḍāvvũ)}{નવડાવવું (navḍāvvũ)}]{\gu{નવડાવવું} (navḍāvvũ) — to bathe someone}

\label{lesson:GU-C488-navdavvu}



\begin{quote}
Before the new one: say the Gujarati for to massage by pressing, then the Gujarati for to cover someone with a blanket.
\end{quote}

\subsection*{You'll want to know: \gu{નવડાવવું}}
\textbf{\gu{નવડાવવું}} — \emph{navḍāvvũ} — \textquotedblleft{}to bathe someone\textquotedblright{}.

\begin{quote}
\gu{માતા} \gu{રોજ} \gu{સવારે} \gu{બાળકને} \gu{નવડાવે} \gu{છે}.
\end{quote}

\begin{quote}
\emph{mātā roj savāre bāḷakne navḍāve che.}
\end{quote}

\begin{quote}
The mother bathes the baby every morning.
\end{quote}

Past, with \gu{તેણે} (he or she …): \gu{તેણે} \gu{બાળકને} \gu{નવડાવ્યું}, \emph{teṇe bāḷakne navḍāvyũ}.

\begin{cousinweb}
\textbf{\gu{નવડાવવું}} (\emph{navḍāvvũ}) is the causative of \textbf{\gu{નાહવું}} (\emph{nāhvũ}), \textquotedblleft{}to bathe\textquotedblright{}, which is inherited from Sanskrit \emph{snāti}, \textquotedblleft{}bathes\textquotedblright{} (CDIAL 13786).
\end{cousinweb}

\begin{grammarlens}[title={What you've built}]
One more word for describing what you lived and felt.
\end{grammarlens}

\subsection*{Your turn}
\begin{itemize}
\raggedright
  \item \emph{Say it:} \emph{navḍāvvũ}
  \item \emph{Say it:} \emph{navḍāvvũ}, once more
  \item \emph{Say it:} \emph{mātā roj savāre bāḷakne navḍāve che.} — The mother bathes the baby every morning.
  \item \emph{Recall:} say the Gujarati for to massage by pressing, then the Gujarati for to cover someone with a blanket, then say \emph{navḍāvvũ} again
\end{itemize}

\begin{tcolorbox}[breakable,colback=teal!4,colframe=teal!35!black,title={Before you move on}]
What does \gu{નવડાવવું} mean? (\textquotedblleft{}To bathe someone\textquotedblright{}.) Say it once more.
\end{tcolorbox}
\section[\texorpdfstring{\gu{ખવડાવવું} (khavḍāvvũ)}{ખવડાવવું (khavḍāvvũ)}]{\gu{ખવડાવવું} (khavḍāvvũ) — to feed}

\label{lesson:GU-C488-khavdavvu}



\begin{quote}
Before the new one: say the Gujarati for to cover someone with a blanket, then the Gujarati for to bathe someone.
\end{quote}

\subsection*{You'll want to know: \gu{ખવડાવવું}}
\textbf{\gu{ખવડાવવું}} — \emph{khavḍāvvũ} — \textquotedblleft{}to feed\textquotedblright{}.

\begin{quote}
\gu{દાદી} \gu{બાળકને} \gu{ચમચીથી} \gu{ખીચડી} \gu{ખવડાવે} \gu{છે}.
\end{quote}

\begin{quote}
\emph{dādī bāḷakne camcīthī khīcḍī khavḍāve che.}
\end{quote}

\begin{quote}
Grandmother feeds the baby khichdi with a spoon.
\end{quote}

Past, with \gu{તેણે} (he or she …): \gu{તેણે} \gu{બાળકને} \gu{ખીચડી} \gu{ખવડાવી}, \emph{teṇe bāḷakne khīcḍī khavḍāvī}.

\begin{cousinweb}
\textbf{\gu{ખવડાવવું}} (\emph{khavḍāvvũ}) is the causative of \textbf{\gu{ખાવું}} (\emph{khāvũ}), \textquotedblleft{}to eat\textquotedblright{}, which goes back to Sanskrit \emph{khādati}, \textquotedblleft{}chews, eats\textquotedblright{} (CDIAL 3865).
\end{cousinweb}

\begin{grammarlens}[title={What you've built}]
One more word for describing what you lived and felt.
\end{grammarlens}

\subsection*{Your turn}
\begin{itemize}
\raggedright
  \item \emph{Say it:} \emph{khavḍāvvũ}
  \item \emph{Say it:} \emph{khavḍāvvũ}, once more
  \item \emph{Say it:} \emph{dādī bāḷakne camcīthī khīcḍī khavḍāve che.} — Grandmother feeds the baby khichdi with a spoon.
  \item \emph{Recall:} say the Gujarati for to cover someone with a blanket, then the Gujarati for to bathe someone, then say \emph{khavḍāvvũ} again
\end{itemize}

\begin{tcolorbox}[breakable,colback=teal!4,colframe=teal!35!black,title={Before you move on}]
What does \gu{ખવડાવવું} mean? (\textquotedblleft{}To feed\textquotedblright{}.) Say it once more.
\end{tcolorbox}
\section[\texorpdfstring{\gu{પહેરાવવું} (paherāvvũ)}{પહેરાવવું (paherāvvũ)}]{\gu{પહેરાવવું} (paherāvvũ) — to dress someone}

\label{lesson:GU-C488-paheravvu}



\begin{quote}
Before the new one: say the Gujarati for to bathe someone, then the Gujarati for to feed.
\end{quote}

\subsection*{You'll want to know: \gu{પહેરાવવું}}
\textbf{\gu{પહેરાવવું}} — \emph{paherāvvũ} — \textquotedblleft{}to dress someone\textquotedblright{}.

\begin{quote}
\gu{મીરાએ} \gu{બાળકને} \gu{નવાં} \gu{કપડાં} \gu{પહેરાવ્યાં}.
\end{quote}

\begin{quote}
\emph{Mīrāe bāḷakne navā̃ kapḍā̃ paherāvyā̃.}
\end{quote}

\begin{quote}
Mira dressed the baby in new clothes.
\end{quote}

Past, with \gu{તેણે} (he or she …): \gu{તેણે} \gu{બાળકને} \gu{ટોપી} \gu{પહેરાવી}, \emph{teṇe bāḷakne ṭopī paherāvī}.

\begin{cousinweb}
\textbf{\gu{પહેરાવવું}} (\emph{paherāvvũ}) is the causative of \textbf{\gu{પહેરવું}} (\emph{pahervũ}), \textquotedblleft{}to wear\textquotedblright{}, which goes back to Sanskrit \emph{paridadhāti}, \textquotedblleft{}puts on clothes\textquotedblright{} (CDIAL 7835).
\end{cousinweb}

\begin{grammarlens}[title={What you've built}]
One more word for describing what you lived and felt.
\end{grammarlens}

\subsection*{Your turn}
\begin{itemize}
\raggedright
  \item \emph{Say it:} \emph{paherāvvũ}
  \item \emph{Say it:} \emph{paherāvvũ}, once more
  \item \emph{Say it:} \emph{Mīrāe bāḷakne navā̃ kapḍā̃ paherāvyā̃.} — Mira dressed the baby in new clothes.
  \item \emph{Recall:} say the Gujarati for to bathe someone, then the Gujarati for to feed, then say \emph{paherāvvũ} again
\end{itemize}

\begin{tcolorbox}[breakable,colback=teal!4,colframe=teal!35!black,title={Before you move on}]
What does \gu{પહેરાવવું} mean? (\textquotedblleft{}To dress someone\textquotedblright{}.) Say it once more.
\end{tcolorbox}
\section[\texorpdfstring{\gu{રમાડવું} (ramāḍvũ)}{રમાડવું (ramāḍvũ)}]{\gu{રમાડવું} (ramāḍvũ) — to play with a child}

\label{lesson:GU-C488-ramadvu}



\begin{quote}
Before the new one: say the Gujarati for to feed, then the Gujarati for to dress someone.
\end{quote}

\subsection*{You'll want to know: \gu{રમાડવું}}
\textbf{\gu{રમાડવું}} — \emph{ramāḍvũ} — \textquotedblleft{}to play with a child\textquotedblright{}.

\begin{quote}
\gu{સાંજે} \gu{દાદા} \gu{બાળકને} \gu{બગીચામાં} \gu{રમાડે} \gu{છે}.
\end{quote}

\begin{quote}
\emph{sā̃je dādā bāḷakne bagīcāmā̃ ramāḍe che.}
\end{quote}

\begin{quote}
In the evening my grandfather plays with the baby in the garden.
\end{quote}

Past, with \gu{તેણે} (he or she …): \gu{તેણે} \gu{બાળકને} \gu{રમાડ્યું}, \emph{teṇe bāḷakne ramāḍyũ}.

\begin{cousinweb}
\textbf{\gu{રમાડવું}} (\emph{ramāḍvũ}) is the causative of \textbf{\gu{રમવું}} (\emph{ramvũ}), \textquotedblleft{}to play\textquotedblright{}, which is inherited from Sanskrit \emph{ramatē}, \textquotedblleft{}takes pleasure in\textquotedblright{} (CDIAL 10628).
\end{cousinweb}

\begin{grammarlens}[title={What you've built}]
One more word for describing what you lived and felt.
\end{grammarlens}

\subsection*{Your turn}
\begin{itemize}
\raggedright
  \item \emph{Say it:} \emph{ramāḍvũ}
  \item \emph{Say it:} \emph{ramāḍvũ}, once more
  \item \emph{Say it:} \emph{sā̃je dādā bāḷakne bagīcāmā̃ ramāḍe che.} — In the evening my grandfather plays with the baby in the garden.
  \item \emph{Recall:} say the Gujarati for to feed, then the Gujarati for to dress someone, then say \emph{ramāḍvũ} again
\end{itemize}

\begin{tcolorbox}[breakable,colback=teal!4,colframe=teal!35!black,title={Before you move on}]
What does \gu{રમાડવું} mean? (\textquotedblleft{}To play with a child\textquotedblright{}.) Say it once more.
\end{tcolorbox}
\section[\texorpdfstring{\gu{પીવડાવવું} (pīvḍāvvũ)}{પીવડાવવું (pīvḍāvvũ)}]{\gu{પીવડાવવું} (pīvḍāvvũ) — to give someone a drink}

\label{lesson:GU-C488-pivdavvu}



\begin{quote}
Before the new one: say the Gujarati for to dress someone, then the Gujarati for to play with a child.
\end{quote}

\subsection*{You'll want to know: \gu{પીવડાવવું}}
\textbf{\gu{પીવડાવવું}} — \emph{pīvḍāvvũ} — \textquotedblleft{}to give someone a drink\textquotedblright{}.

\begin{quote}
\gu{નર્સે} \gu{દર્દીને} \gu{પાણી} \gu{પીવડાવ્યું}.
\end{quote}

\begin{quote}
\emph{narse dardīne pāṇī pīvḍāvyũ.}
\end{quote}

\begin{quote}
The nurse gave the patient some water.
\end{quote}

Past, with \gu{તેણે} (he or she …): \gu{તેણે} \gu{બાળકને} \gu{દૂધ} \gu{પીવડાવ્યું}, \emph{teṇe bāḷakne dūdh pīvḍāvyũ}.

\begin{cousinweb}
\textbf{\gu{પીવડાવવું}} (\emph{pīvḍāvvũ}) is the causative of \textbf{\gu{પીવું}} (\emph{pīvũ}), \textquotedblleft{}to drink\textquotedblright{}, which is inherited from Sanskrit \emph{pibati}, \textquotedblleft{}drinks\textquotedblright{} (CDIAL 8209).
\end{cousinweb}

\begin{grammarlens}[title={What you've built}]
One more word for describing what you lived and felt.
\end{grammarlens}

\subsection*{Your turn}
\begin{itemize}
\raggedright
  \item \emph{Say it:} \emph{pīvḍāvvũ}
  \item \emph{Say it:} \emph{pīvḍāvvũ}, once more
  \item \emph{Say it:} \emph{narse dardīne pāṇī pīvḍāvyũ.} — The nurse gave the patient some water.
  \item \emph{Recall:} say the Gujarati for to dress someone, then the Gujarati for to play with a child, then say \emph{pīvḍāvvũ} again
\end{itemize}

\begin{tcolorbox}[breakable,colback=teal!4,colframe=teal!35!black,title={Before you move on}]
What does \gu{પીવડાવવું} mean? (\textquotedblleft{}To give someone a drink\textquotedblright{}.) Say it once more.
\end{tcolorbox}
